\documentclass[11pt]{article}
% ======================================================================
%  examstyle.tex — EPFL-style exam layout for the weekly Time Series exams
%  Shared by every Exam_Week_NN(.tex / _Solutions.tex). Compiles with pdflatex.
% ======================================================================
\usepackage[utf8]{inputenc}
\usepackage[T1]{fontenc}
\usepackage[english]{babel}
\usepackage[a4paper,margin=2.4cm]{geometry}
\usepackage{amsmath,amssymb,amsthm,mathtools}
\usepackage{bm}
\usepackage{enumitem}
\usepackage{array}

\setlength{\parindent}{0pt}
\setlength{\parskip}{0.5em}
\emergencystretch=3em
\allowdisplaybreaks[1]

% ---------- math macros (same names as the rest of the project) ----------
\newcommand{\E}{\mathbb{E}}
\DeclareMathOperator{\Var}{Var}
\DeclareMathOperator{\Cov}{Cov}
\DeclareMathOperator{\Corr}{Corr}
\DeclareMathOperator*{\argmin}{arg\,min}
\DeclareMathOperator{\MSE}{MSE}
\DeclareMathOperator{\trace}{tr}
\newcommand{\Reals}{\mathbb{R}}
\newcommand{\Ints}{\mathbb{Z}}
\newcommand{\Comp}{\mathbb{C}}
\newcommand{\Nat}{\mathbb{N}}
\newcommand{\Prob}{\mathbb{P}}
\newcommand{\Tr}{^{\mathsf{T}}}
\newcommand{\conj}[1]{\overline{#1}}
\newcommand{\abs}[1]{\left\lvert #1 \right\rvert}
\newcommand{\norm}[1]{\left\lVert #1 \right\rVert}
\newcommand{\ind}{\mathbf{1}}
\newcommand{\eps}{\varepsilon}
\newcommand{\diff}{\nabla}
\newcommand{\acvs}{\gamma}
\newcommand{\acf}{\rho}
\newcommand{\pacf}{\alpha}
\newcommand{\sdf}{\mathcal{S}}
\newcommand{\Bsh}{B}
\newcommand{\iid}{\overset{\text{iid}}{\sim}}

\setlist[enumerate]{leftmargin=2em,topsep=2pt,itemsep=4pt}
\setlist[itemize]{leftmargin=1.6em,topsep=2pt,itemsep=2pt}

\newcounter{exo}

% ---------- exam cover header :  \examheader{exam title}{duration} ----------
\newcommand{\examheader}[2]{%
  \thispagestyle{empty}
  \begin{center}
    {\Large\bfseries Time Series}\\[3pt]
    Spring Semester 2026, EPFL\\
    \'Ecole Polytechnique F\'ed\'erale de Lausanne\\
    Prof.\ Dr.\ Sofia C.\ Olhede\\[7pt]
    \hrule height 0.8pt
    \vspace{9pt}
    {\large\bfseries Practice Exam --- #1}\\[4pt]
    {\normalsize Duration: #2 \quad$\cdot$\quad No calculator \quad$\cdot$\quad Answer on paper}\\
    \vspace{8pt}
    \hrule height 0.8pt
  \end{center}
  \vspace{14pt}
  \begin{tabular}{@{}p{8.2cm}@{}p{8cm}@{}}
    Name:\enspace\hrulefill & Surname:\enspace\hrulefill
  \end{tabular}\\[14pt]
  SCIPER (Student Card Number):\enspace\hrulefill
  \vspace{16pt}
}

% ---------- points table (fixed: 4 problems) ----------
\newcommand{\pointstable}{%
  \begin{center}
  \renewcommand{\arraystretch}{1.5}
  \begin{tabular}{|p{5cm}|p{4cm}|}
    \hline
    \textbf{Exercise} & \textbf{Points}\\\hline
    1 & \\\hline
    2 & \\\hline
    3 & \\\hline
    4 & \\\hline
    \textbf{Total} & \\\hline
  \end{tabular}
  \end{center}
}

% ---------- problem heading (exam: one problem per page) ----------
\newcommand{\problem}[1]{%
  \clearpage
  \stepcounter{exo}%
  \begin{center}\textbf{\large Problem \theexo\ (#1 points)}\end{center}
  \vspace{4pt}
}
% ---------- answer space: fill the statement page + one blank answer page ----------
% Put \answerpage at the END of each problem so every exercise gets its own page(s).
\newcommand{\answerpage}{\par\vfill\newpage\null\newpage}

% ---------- solutions document ----------
\newcommand{\solheader}[1]{%
  \begin{center}
    {\Large\bfseries Time Series --- Solutions}\\[3pt]
    {\large Practice Exam --- #1}\\[2pt]
    EPFL, MATH-342 \quad$\cdot$\quad Prof.\ S.\ C.\ Olhede\\[5pt]
    \hrule height 0.8pt
  \end{center}
  \vspace{10pt}
}
\newcommand{\solproblem}[1]{%
  \bigskip
  \stepcounter{exo}%
  {\large\bfseries Problem \theexo\ (#1 points)}\par\nobreak\vspace{2pt}
}
% Rappel de l'énoncé avant la solution (dans les corrigés)
\newcommand{\statementstart}{\par\vspace{1pt}{\bfseries\itshape Statement.}\par\nobreak\begingroup\small\itshape}
\newcommand{\statementend}{\par\endgroup\vspace{5pt}\hrule\vspace{6pt}{\bfseries Solution.}\par\nobreak\vspace{2pt}}

\begin{document}

\examheader{Week 4: Parametric Estimation --- Yule--Walker \& Least Squares}{60 minutes}
\pointstable
\begin{center}\itshape Justify every answer. Throughout, $\eps_t$ is a mean-zero white noise of variance $\sigma^2$.\end{center}

% ---------------------------------------------------------------
\problem{5}
Let $\{X_t\}$ be a mean-zero causal $\mathrm{AR}(2)$ process
$X_t=\phi_1X_{t-1}+\phi_2X_{t-2}+\eps_t$.

\textbf{(a)} Starting from the model equation, derive the \emph{Yule--Walker equations}: multiply by $X_{t-k}$, take expectations, and explain precisely where \emph{causality} is used to kill the noise cross-term. State the resulting linear system $\boldsymbol\acvs=\Gamma\boldsymbol\phi$ and the variance identity for $\sigma^2$. \hfill(2 pts)

\textbf{(b)} A series gives the sample autocovariances $\hat\acvs_0=5$, $\hat\acvs_1=3$, $\hat\acvs_2=1$. Compute the Yule--Walker estimates $\hat\phi_1,\hat\phi_2$ and $\hat\sigma^2$. \hfill(2 pts)

\textbf{(c)} Verify that the fitted model of part (b) is causal (stationary) by locating the roots of its AR polynomial. \hfill(1 pt)
\answerpage

% ---------------------------------------------------------------
\problem{5}
A series of length $n=150$ is modelled as a mean-zero $\mathrm{AR}(3)$ process
$X_t=\phi_1X_{t-1}+\phi_2X_{t-2}+\phi_3X_{t-3}+\eps_t$. The estimated
autocorrelations are
\[
\hat\acf_1=0.8,\qquad \hat\acf_2=0.6,\qquad \hat\acf_3=0.4,
\]
and the sample variance is $\hat\acvs_0=10$.

\textbf{(a)} Explain why the Yule--Walker fit can be written entirely with the \emph{correlation} matrix, $\hat{\boldsymbol\phi}=\bar\Gamma_3^{-1}\hat{\boldsymbol\acf}_3$, i.e.\ why $\hat\acvs_0$ cancels in the system for $\boldsymbol\phi$. Write out the $3\times3$ system explicitly. \hfill(1.5 pts)

\textbf{(b)} Solve the system for $\hat\phi_1,\hat\phi_2,\hat\phi_3$ (e.g.\ by Gaussian elimination) and verify your solution by substituting it back into the three equations. \hfill(2.5 pts)

\textbf{(c)} Compute the innovation-variance estimate $\hat\sigma^2$ using $\hat\acvs_0=10$. \hfill(1 pt)
\answerpage

% ---------------------------------------------------------------
\problem{5}
This problem treats the two faces of least squares.

\textbf{(a)} \emph{(MA via residual reconstruction.)} A mean-zero $\mathrm{MA}(1)$ process $X_t=\eps_t-\theta\eps_{t-1}$ is observed at $X_1=0$, $X_2=2$, $X_3=-1$. Reconstruct the innovations $\hat\eps_1,\hat\eps_2,\hat\eps_3$ as functions of $\theta$ (state your initialisation), write the least-squares criterion $S(\theta)$, and find $\hat\theta$. Is the fit invertible? \hfill(2.5 pts)

\textbf{(b)} \emph{(Forward least squares for $\mathrm{AR}(2)$.)} For the five observations $x_1,\dots,x_5=1,\,2,\,0,\,-1,\,-2$ (already mean-zero), write the forward design $\mathbf X_F=F\boldsymbol\phi+\boldsymbol\eps$ for an $\mathrm{AR}(2)$ fit, form the normal equations $F\Tr F\,\hat{\boldsymbol\phi}_F=F\Tr\mathbf X_F$, and solve for $\hat{\boldsymbol\phi}_F$. Then report $\hat\sigma_F^2$ using the correct degrees of freedom. \hfill(2 pts)

\textbf{(c)} State, in one sentence each, (i) why the \emph{backward} least-squares estimator $\hat{\boldsymbol\phi}_B$ estimates the \emph{same} coefficients $\boldsymbol\phi$, and (ii) whether a least-squares AR fit is guaranteed to be stationary. \hfill(0.5 pt)
\answerpage

% ---------------------------------------------------------------
\problem{5}
\textbf{(Revision of Week 3.)} Consider the $\mathrm{ARMA}(2,1)$ process
\[
\Big(1-\tfrac34\Bsh+\tfrac18\Bsh^2\Big)X_t=\Big(1+\tfrac13\Bsh\Big)\eps_t .
\]

\textbf{(a)} Factor the AR polynomial $\Phi(z)=1-\tfrac34 z+\tfrac18 z^2$, locate its roots, and confirm the process is causal. Then check invertibility from $\Theta(z)=1+\tfrac13 z$. \hfill(2 pts)

\textbf{(b)} Using a \emph{partial-fraction} expansion of $G(z)=\Theta(z)/\Phi(z)$, find a closed form for the $\mathrm{MA}(\infty)$ weights $\psi_j$ in $X_t=\sum_{j\ge0}\psi_j\eps_{t-j}$, and report $\psi_0,\psi_1,\psi_2$. \hfill(3 pts)
\answerpage

\end{document}
